% 第5章《连续函数》题目解答
% 对应 chapters/ch05.tex；按原章顺序排列。

\chapter*{第5章\quad 连续函数习题解答}
\addcontentsline{toc}{chapter}{第5章\quad 连续函数习题解答}
\subsection*{5.1 连续性概念}

\subsubsection*{5.1.2 思考题}

\paragraph{思考题 1}
\begin{xsolution}
若 $f$ 在 $a$ 连续，则由连续函数的四则运算以及绝对值函数的连续性，$f^2$ 与 $|f|$ 都在 $a$ 连续。

反之均不成立。例如令
\[
 f(x)=\begin{cases}
  1,&x\ge a,\\
  -1,&x<a.
 \end{cases}
\]
则 $f$ 在 $a$ 不连续，而 $f^2\equiv1$、$|f|\equiv1$ 均在 $a$ 连续。因此仅由 $f^2$ 或 $|f|$ 在 $a$ 连续不能推出 $f$ 在 $a$ 连续。
\end{xsolution}

\paragraph{思考题 2}
\begin{xsolution}
不能由“$f,g$ 在 $a$ 都不连续”唯一判断 $f+g$ 或 $fg$ 是否连续。

例如取 Dirichlet 函数 $D$，令 $f=D,g=1-D$，则 $f,g$ 处处不连续，但
\[
 f+g\equiv1,\qquad fg\equiv0
\]
处处连续。另一方面，若取 $f=g=D$，则 $f+g=2D$、$fg=D$ 都处处不连续。因此和与积既可能连续，也可能不连续。
\end{xsolution}

\paragraph{思考题 3}
\begin{xsolution}
任取 $x_0\in(a,b)$。可取 $c,d$ 使
\[
 a<c<x_0<d<b.
\]
由题设，$f$ 在闭区间 $[c,d]$ 上连续，特别地在其内点 $x_0$ 连续。由于 $x_0$ 任意，$f$ 在 $(a,b)$ 上处处连续，即 $f\in C(a,b)$。
\end{xsolution}

\paragraph{思考题 4}
\begin{xsolution}
\begin{enumerate}[label=(\arabic*)]
 \item 当 $x\to0^-$ 时，$1/x\to-\infty$，故 $\ee^{1/x}\to0$；当 $x\to0^+$ 时，$x^2\to0$，且 $f(0)=0$。所以 $f$ 在 $0$ 连续。其余各点由初等函数连续性知均连续。因此 $f$ 在 $\RR$ 上连续。
 \item 除 $0$ 外均连续。在 $0$ 处
 \[
  \lim_{x\to0^-}\ee^{1/x}=0,\qquad
  \lim_{x\to0^+}\ee^{1/x}=+\infty,
 \]
 两个单侧极限不能同时作为有限实数存在，故 $0$ 为第二类间断点。
 \item 对 $x\ne0$，
 \[
  \frac{\ln(1+x^2)}x
  =x\,\frac{\ln(1+x^2)}{x^2}\longrightarrow0
  \qquad(x\to0),
 \]
 因为 $\ln(1+u)/u\to1$。而 $f(0)=2$，所以 $0$ 为可去间断点；在 $(-1,0)\cup(0,1)$ 上函数连续。
\end{enumerate}
\end{xsolution}

\paragraph{思考题 5}
\begin{xsolution}
\begin{enumerate}[label=(\arabic*)]
 \item $\operatorname{sgn}x$ 只在 $x=0$ 间断，且
 \[
  f(0^-)=-1,\qquad f(0^+)=1,
 \]
 故 $0$ 为跳跃间断点。

 \item $g(x)=x-[x]$ 为小数部分函数。对每个 $m\in\ZZ$，
 \[
  g(m^-)=1,\qquad g(m^+)=g(m)=0,
 \]
 故所有整数都是跳跃间断点；其余点连续。

 \item 由 $g(x)\in[0,1)$，且 $g(x)=0$ 当且仅当 $x\in\ZZ$，有
 \[
  f(g(x))=\begin{cases}0,&x\in\ZZ,\\1,&x\notin\ZZ.
  \end{cases}
 \]
 对任意整数 $m$，$\lim_{x\to m}f(g(x))=1$ 而函数值为 $0$，故每个整数都是可去间断点；其余点连续。

 \item 因为 $f(x)\in\{-1,0,1\}\subset\ZZ$，而 $g(k)=0$ 对每个整数 $k$ 成立，所以
 \[
  g(f(x))\equiv0,
 \]
 因而处处连续。

 \item 原式定义域排除 $x=-1,0,1$。在定义域内可化为
 \[
  h(x)=\frac{x-1}{x+1}.
 \]
 因此 $x=0,1$ 均为可去间断点；在 $x=-1$ 处函数无界，故为第二类间断点。

 \item 定义域为
 \[
  (-\infty,0)\cup(0,1],
 \]
 因为当 $x>1$ 时 $[1/x]=0$。在每个使 $1/x$ 不越过整数的开区间内函数为常数。对 $n\in\NN_+$，点 $x=-1/n$ 的左右函数值分别趋于 $-1/n$ 与 $-1/(n+1)$，故都是跳跃点；对 $n\ge2$，点 $x=1/n$ 的左右函数值分别为 $1/n$ 与 $1/(n-1)$，故也是跳跃点。点 $x=1$ 按定义域的左侧连续。因此全部间断点为
 \[
  \left\{-\frac1n:n\in\NN_+\right\}
  \cup
  \left\{\frac1n:n\ge2\right\},
 \]
 且均为跳跃间断点。
\end{enumerate}
\end{xsolution}

\subsubsection*{5.1.4 练习题}

\paragraph{练习题 1}
\begin{xsolution}
\begin{enumerate}[label=(\arabic*)]
 \item 连续性第一定义的否定可写成：存在 $\varepsilon_0>0$，使对每个 $\delta>0$，总存在定义域中的 $x$ 满足
 \[
  0<|x-a|<\delta,\qquad |f(x)-f(a)|\ge\varepsilon_0.
 \]
 连续性第二定义的否定可写成：存在数列 $\{x_n\}$，$x_n\to a$，但 $f(x_n)\not\to f(a)$。等价地，可在取子列后写成存在 $\varepsilon_0>0$ 及 $x_n\to a$，使
 \[
  |f(x_n)-f(a)|\ge\varepsilon_0,\qquad n\in\NN_+.
 \]

 \item 先设第一定义成立。若 $x_n\to a$，则对任意 $\varepsilon>0$，由连续性可取 $\delta>0$，使 $|x-a|<\delta$ 时 $|f(x)-f(a)|<\varepsilon$。当 $n$ 充分大时 $|x_n-a|<\delta$，故 $|f(x_n)-f(a)|<\varepsilon$，于是 $f(x_n)\to f(a)$。

 反过来，若第一定义不成立，则由 (1) 存在 $\varepsilon_0>0$，对每个 $n$ 可取 $x_n$ 使
 \[
  0<|x_n-a|<\frac1n,
  \qquad |f(x_n)-f(a)|\ge\varepsilon_0.
 \]
 于是 $x_n\to a$，但 $f(x_n)\not\to f(a)$，与第二定义矛盾。因此两定义等价。
\end{enumerate}
\end{xsolution}

\paragraph{练习题 2}
\begin{xsolution}
\begin{enumerate}[label=(\arabic*)]
 \item $[x]$ 在每个非整数点连续；对 $m\in\ZZ$，有
 \[
  [x]_{x\to m^-}=m-1,\qquad [m]=[x]_{x\to m^+}=m,
 \]
 所以所有整数均为跳跃间断点。

 \item 当 $x\to0$ 时
 \[
  \frac{\ee^{\sin x}-1}{x}
  =\frac{\ee^{\sin x}-1}{\sin x}\frac{\sin x}{x}\longrightarrow1=f(0).
 \]
 其余点处由初等函数连续性可知连续，故该函数处处连续。

 \item 若 $x\notin\ZZ$，则 $[x]+[-x]=-1$；若 $x\in\ZZ$，则该值为 $0$。所以每个整数处的极限均为 $-1$，函数值为 $0$，故所有整数都是可去间断点，其余点连续。

 \item 当 $x\ne0,\pm1$ 时
 \[
  \frac{x^2-x}{|x|(x^2-1)}
  =\frac{x}{|x|(x+1)}
  =\frac{\operatorname{sgn}x}{x+1}.
 \]
 因而在 $x=0$，左右极限分别为 $-1,1$，是跳跃点；在 $x=1$，极限为 $1/2$ 而函数值为 $1$，是可去间断点；在 $x=-1$，两侧分别无界，是第二类间断点。
\end{enumerate}
\end{xsolution}

\paragraph{练习题 3}
\begin{xsolution}
数列 $\{x_n\}\subset[a,b]$ 有收敛子列 $x_{n_k}\to\xi\in[a,b]$。由 $f$ 在 $\xi$ 连续，
\[
 f(x_{n_k})\to f(\xi).
\]
另一方面原数列满足 $f(x_n)\to A$，故其任意子列也趋于 $A$。由极限唯一性，$f(\xi)=A$。
\end{xsolution}

\paragraph{练习题 4}
\begin{xsolution}
若 $0<x<1$，则反复使用 $f(x)=f(x^2)$ 得
\[
 f(x)=f(x^{2^n}),
\]
而 $x^{2^n}\to0$，由 $f$ 在 $0$ 连续知 $f(x)=f(0)$。

若 $x>1$，令 $y_n=x^{1/2^n}$。由 $f(y)=f(y^2)$ 逐次反推得 $f(x)=f(y_n)$，而 $y_n\to1$，故由 $f$ 在 $1$ 连续，$f(x)=f(1)$。

又令 $x\to1^-$，前一结论给出 $f(x)=f(0)$，连续性于是推出 $f(1)=f(0)$。最后，对 $x<0$，有 $x^2>0$，故
\[
 f(x)=f(x^2)=f(0).
\]
所以 $f$ 在 $\RR$ 上为常值函数。
\end{xsolution}

\paragraph{练习题 5}
\begin{xsolution}
利用恒等式
\[
 \max\{f,g\}=\frac{f+g+|f-g|}{2},
 \qquad
 \min\{f,g\}=\frac{f+g-|f-g|}{2}.
\]
因为 $f,g$ 连续，所以 $f\pm g$ 连续，进而 $|f-g|$ 连续。因此上面两个函数均连续，即属于 $C(I)$。
\end{xsolution}

\paragraph{练习题 6}
\begin{xsolution}
记
\[
 M(x)=\max\{f_1(x),f_2(x),f_3(x)\},
 \qquad
 m(x)=\min\{f_1(x),f_2(x),f_3(x)\}.
\]
由上一题反复应用于两个函数可知 $M,m\in C[a,b]$。三个数的中间值等于三数之和减去最大值和最小值，所以
\[
 f(x)=f_1(x)+f_2(x)+f_3(x)-M(x)-m(x).
\]
右端为连续函数之和差，故 $f\in C[a,b]$。
\end{xsolution}

\paragraph{练习题 7}
\begin{xsolution}
若 $f$ 连续，则
\[
 f_n(x)=\max\{-n,\min\{f(x),n\}\}
\]
由连续函数的最大、最小运算仍连续，故每个 $f_n$ 连续。

反之，设每个 $f_n$ 连续。固定 $x_0$，取正整数 $n>|f(x_0)|+1$。此时 $f_n(x_0)=f(x_0)$。由 $f_n$ 在 $x_0$ 连续，存在 $x_0$ 的一个邻域，使
\[
 |f_n(x)-f(x_0)|<1.
\]
于是该邻域内 $|f_n(x)|<n$。若某点有 $f(x)\ge n$ 或 $f(x)\le-n$，则 $f_n(x)$ 必等于 $n$ 或 $-n$，与 $|f_n(x)|<n$ 矛盾。因此在此邻域内 $f_n(x)=f(x)$。故 $f$ 在 $x_0$ 附近与连续函数 $f_n$ 相同，从而在 $x_0$ 连续。$x_0$ 任意，所以 $f$ 连续。
\end{xsolution}

\paragraph{练习题 8}
\begin{xsolution}
设 Dirichlet 函数
\[
 D(x)=\begin{cases}1,&x\in\QQ,\\0,&x\notin\QQ.
 \end{cases}
\]
任意点 $a$ 的任意邻域中都有有理数和无理数。取有理数列 $r_n\to a$ 与无理数列 $s_n\to a$，则
\[
 D(r_n)=1,\qquad D(s_n)=0.
\]
所以 $D(x)$ 在 $a$ 处极限不存在。故 $D$ 处处不连续，并且每个间断点都是第二类间断点。
\end{xsolution}

\paragraph{练习题 9}
\begin{xsolution}
给定指定点 $a$，令
\[
 f(x)=|x-a|D(x),
\]
其中 $D$ 为 Dirichlet 函数。在 $x=a$ 处，
\[
 |f(x)|\le|x-a|\to0=f(a),
\]
故 $f$ 在 $a$ 连续。

若 $x_0\ne a$，沿有理数列趋于 $x_0$ 时函数值趋于 $|x_0-a|>0$，沿无理数列趋于 $x_0$ 时函数值恒为 $0$，故极限不存在。因此除 $a$ 外处处不连续。
\end{xsolution}

\paragraph{练习题 10}
\begin{xsolution}
令 $y=0$，得 $f(x)=f(x)f(0)$。若存在 $x$ 使 $f(x)\ne0$，则 $f(0)=1$；若不存在，则 $f\equiv0$。

以下设 $f\not\equiv0$。由
\[
 f(x)=f(x/2)^2\ge0.
\]
若存在 $x_0$ 使 $f(x_0)=0$，则
\[
 1=f(0)=f(x_0)f(-x_0)=0,
\]
矛盾，故 $f(x)>0$ 对所有 $x$ 成立。定义
\[
 g(x)=\ln f(x).
\]
则 $g$ 连续且
\[
 g(x+y)=g(x)+g(y).
\]
连续的加法函数必为线性函数，故 $g(x)=cx$，其中 $c=g(1)$。于是
\[
 f(x)=\ee^{cx}=a^x,
 \qquad a=\ee^c=f(1)>0.
\]
反之 $f\equiv0$ 与任意 $a>0$ 的 $f(x)=a^x$ 都满足原方程。
\end{xsolution}

\paragraph{练习题 11}
\begin{xsolution}
令
\[
 g(x)=f(x)-f(0).
\]
则 $g$ 连续、$g(0)=0$，且同样满足中点等式
\[
 g\left(\frac{x+y}{2}\right)=\frac{g(x)+g(y)}2.
\]
取 $y=0$ 得 $g(x/2)=g(x)/2$，等价于 $g(2u)=2g(u)$。于是
\[
 g(x+y)=2g\left(\frac{x+y}{2}\right)=g(x)+g(y).
\]
所以 $g$ 是连续加法函数，故 $g(x)=g(1)x$。因此
\[
 f(x)=[f(1)-f(0)]x+f(0).
\]
\end{xsolution}

\paragraph{练习题 12}
\begin{xsolution}
若 $f$ 在 $a$ 连续，则给定 $\varepsilon>0$，可取 $\delta>0$，使当 $x\in O_\delta(a)$ 时
\[
 |f(x)-f(a)|<\frac\varepsilon2.
\]
于是任意 $x,y\in O_\delta(a)$ 有
\[
 |f(x)-f(y)|\le|f(x)-f(a)|+|f(y)-f(a)|<\varepsilon.
\]
所以 $\omega_f(a,\delta)\le\varepsilon$，令 $\delta\to0^+$ 得 $\omega_f(a)=0$。

反之，若 $\omega_f(a)=0$，则对任意 $\varepsilon>0$，存在 $\delta>0$ 使
\[
 \omega_f(a,\delta)<\varepsilon.
\]
因为 $a,x\in O_\delta(a)$ 时
\[
 |f(x)-f(a)|\le
 \sup_{O_\delta(a)}f-\inf_{O_\delta(a)}f
 =\omega_f(a,\delta)<\varepsilon,
\]
故 $f$ 在 $a$ 连续。
\end{xsolution}

\subsection*{5.2 零点存在定理与介值定理}

\paragraph{5.2.1 思考题}
\begin{xsolution}
可以构造一个在每个非退化区间上都取遍 $\RR$ 的函数。把 $\RR$ 按等价关系
\[
 x\sim y\quad\Longleftrightarrow\quad x-y\in\QQ
\]
分成陪集。每个陪集 $x+\QQ$ 都在 $\RR$ 中稠密，而且陪集全体的势与 $\RR$ 相同。取一个双射
\[
 \Phi\colon \RR/\QQ\to\RR,
\]
并定义 $f(x)=\Phi(x+\QQ)$。

任意非退化区间都与每个陪集相交，所以它在 $f$ 下的像就是整个 $\RR$。因此 $f$ 具有介值性质。另一方面，任意点的任意邻域内 $f$ 都取遍 $\RR$，故 $f$ 在每一点都不连续。
\end{xsolution}

\subsubsection*{5.2.3 练习题}

\paragraph{练习题 1}
\begin{xsolution}
设
\[
 m=\min_{1\le i\le n}f(x_i),\qquad
 M=\max_{1\le i\le n}f(x_i).
\]
算术平均数
\[
 A=\frac1n\sum_{i=1}^n f(x_i)
\]
满足 $m\le A\le M$。取达到 $m,M$ 的两个点，它们都属于 $[x_1,x_n]$。由连续函数的介值定理，$f$ 在连接这两点的区间上取遍 $[m,M]$，故存在 $\xi\in[x_1,x_n]$ 使 $f(\xi)=A$。
\end{xsolution}

\paragraph{练习题 2}
\begin{xsolution}
用参数 $\theta$ 表示圆周上的点，令 $F(\theta)$ 表示函数在该点的值。所谓圆周上的连续，是指 $F$ 为 $2\pi$ 周期连续函数。定义
\[
 h(\theta)=F(\theta)-F(\theta+\pi).
\]
则 $h$ 连续，且
\[
 h(\theta+\pi)=-h(\theta).
\]
若 $h(\theta_0)=0$，则以 $\theta_0$ 与 $\theta_0+\pi$ 为端点的直径即满足要求。否则 $h(\theta_0)$ 与 $h(\theta_0+\pi)$ 异号，由零点存在定理，在二者之间必有 $\xi$ 使 $h(\xi)=0$。故必存在所求直径。
\end{xsolution}

\paragraph{练习题 3}
\begin{xsolution}
题设条件
\[
 \sum_{k=1}^{n-1}|a_k|<a_n
\]
在一般情形下不足以推出结论，因为它没有约束常数项 $a_0$。例如取 $n=2$，令
\[
 a_2=1,\qquad a_1=0,\qquad a_0=2.
\]
则 $|a_1|<a_2$ 成立，但
\[
 C_2(x)=2+\cos2x\ge1>0,
\]
在 $[0,2\pi)$ 中没有零点。因此按题设条件，命题不成立。

若把条件加强为
\[
 \sum_{k=0}^{n-1}|a_k|<|a_n|,
\]
则结论成立。令 $x_j=j\pi/n$，$j=0,1,\ldots,2n$。由于
\[
 a_n\cos(nx_j)=a_n(-1)^j
\]
且其绝对值严格大于其余各项绝对值之和，所以 $C_n(x_j)$ 的符号随 $j$ 交替。于是对每个 $j=0,1,\ldots,2n-1$，连续函数 $C_n$ 在 $(x_j,x_{j+1})$ 内至少有一个零点，从而在 $[0,2\pi)$ 中至少有 $2n$ 个零点。
\end{xsolution}

\paragraph{练习题 4}
\begin{xsolution}
记
\[
 F(x)=\frac{a_1}{x-b_1}+\frac{a_2}{x-b_2}+\frac{a_3}{x-b_3}.
\]
在每个不含 $b_i$ 的区间内 $F$ 连续。若 $x<y$ 且二者同处于 $(b_1,b_2)$ 或 $(b_2,b_3)$，则对每个 $i$，函数 $1/(x-b_i)$ 随 $x$ 严格减少；因 $a_i>0$，故 $F$ 在这两个区间上都严格减少。

在 $(b_1,b_2)$ 上，
\[
 \lim_{x\to b_1^+}F(x)=+\infty,
 \qquad
 \lim_{x\to b_2^-}F(x)=-\infty.
\]
由介值定理至少有一根，由严格单调性至多有一根，故恰有一根。同理，在 $(b_2,b_3)$ 上
\[
 \lim_{x\to b_2^+}F(x)=+\infty,
 \qquad
 \lim_{x\to b_3^-}F(x)=-\infty,
\]
也恰有一根。
\end{xsolution}

\paragraph{练习题 5}
\begin{xsolution}
\begin{enumerate}[label=(\arabic*)]
 \item 奇数次首一多项式满足
 \[
  \lim_{x\to+\infty}f(x)=+\infty,
  \qquad
  \lim_{x\to-\infty}f(x)=-\infty.
 \]
 因此可取 $A<0<B$ 使 $f(A)<0<f(B)$。由零点存在定理，$f$ 至少有一个实根。

 \item 偶数次多项式可以没有实根，例如 $x^2+1=0$。若常数项 $a_{2n}<0$，则 $f(0)<0$，而
 \[
  \lim_{x\to\pm\infty}f(x)=+\infty.
 \]
 因而可取 $A<0<B$ 使 $f(A)>0,f(B)>0$。分别在 $[A,0]$ 和 $[0,B]$ 上应用零点存在定理，各得到一个实根；两根一负一正，故互异。
\end{enumerate}
\end{xsolution}

\paragraph{练习题 6}
\begin{xsolution}
令
\[
 F(x)=x^{17}+\frac{215}{1+\cos^2 3x}-18.
\]
$F$ 在 $\RR$ 上连续。因
\[
 F(0)=\frac{215}{2}-18>0,
\]
而
\[
 F(-2)\le -2^{17}+215-18<0,
\]
故由零点存在定理，$F$ 在 $(-2,0)$ 内有零点，即原方程必有实根。
\end{xsolution}

\paragraph{练习题 7}
\begin{xsolution}
连续函数把区间映成区间，所以 $f([a,b])$ 是 $\RR$ 中的区间。另一方面题设说该值域只含有理数。若值域中有两个不同的数 $p<q$，则区间性质要求它包含 $p,q$ 之间的所有实数，其中必有无理数，矛盾。因此值域只能退化为一个点，故 $f$ 是常值函数。
\end{xsolution}

\paragraph{练习题 8}
\begin{xsolution}
先证明一般事实：区间上的连续一一映射必严格单调。否则存在 $x_1<x_2<x_3$，使 $f(x_2)$ 不严格位于 $f(x_1),f(x_3)$ 之间。若例如
\[
 f(x_2)>\max\{f(x_1),f(x_3)\},
\]
取 $c$ 严格介于 $\max\{f(x_1),f(x_3)\}$ 与 $f(x_2)$ 之间，则由介值定理，在 $(x_1,x_2)$ 与 $(x_2,x_3)$ 各有一点取值 $c$，与一一性矛盾。其余情形同理。因此 $f$ 严格单调。

若 $f(a)<f(b)$，严格单调不可能为减少，只能严格增加；若 $f(a)>f(b)$，则只能严格减少。这就得到 (1)、(2)。
\end{xsolution}

\paragraph{练习题 9}
\begin{xsolution}
给定 $c\in(A,B)$。由
\[
 \lim_{x\to-\infty}f(x)=A<c,
 \qquad
 \lim_{x\to+\infty}f(x)=B>c,
\]
可取 $u<v$，使
\[
 f(u)<c<f(v).
\]
函数 $f-c$ 在 $[u,v]$ 连续且两端异号，由零点存在定理存在 $\xi\in(u,v)$ 使 $f(\xi)=c$。
\end{xsolution}

\paragraph{练习题 10}
\begin{xsolution}
不妨设 $A<\eta<B$。由 $f(x_n)\to A$、$f(y_n)\to B$，当 $n$ 充分大时有
\[
 f(x_n)<\eta<f(y_n).
\]
对连接 $x_n,y_n$ 的闭区间应用介值定理，可取位于二者之间的 $z_n$，使
\[
 f(z_n)=\eta.
\]
对有限个较小的 $n$ 任意补取 $z_n\in[a,b)$ 即可。因为 $x_n\to b,y_n\to b$，而 $z_n$ 位于 $x_n,y_n$ 之间，所以夹逼得 $z_n\to b$；同时 $f(z_n)=\eta$ 对充分大的 $n$ 恒成立，故 $f(z_n)\to\eta$。
\end{xsolution}

\subsection*{5.3 有界性定理与最值定理}

\subsubsection*{5.3.3 练习题}

\paragraph{练习题 1}
\begin{xsolution}
任取 $x_0\in[a,b]$。若 $x_0\in(a,b)$，题设保证 $f(x_0^-),f(x_0^+)$ 都是有限数；由单侧极限的局部有界性，存在 $x_0$ 的一个邻域，使 $f$ 在去掉 $x_0$ 后有界，再把有限的函数值 $f(x_0)$ 加入即可知 $f$ 在该邻域内有界。端点处只需使用相应的单侧极限，同样得到局部有界性。

于是 $[a,b]$ 的每一点都有一个邻域，在该邻域与 $[a,b]$ 的交上 $f$ 有界。这些邻域构成 $[a,b]$ 的开覆盖。由有限覆盖定理可取有限子覆盖。有限多个局部界的最大值给出整个 $[a,b]$ 上的统一界，故 $f$ 有界。
\end{xsolution}

\paragraph{练习题 2}
\begin{xsolution}
设
\[
 \lim_{x\to+\infty}f(x)=A\in\RR.
\]
取 $M>a$，使 $x>M$ 时 $|f(x)-A|<1$，于是尾部 $[M,+\infty)$ 上
\[
 |f(x)|\le |A|+1.
\]
在闭区间 $[a,M]$ 上，$f$ 连续，故由有界性定理有界。合并两部分即可知 $f$ 在 $[a,+\infty)$ 上有界。
\end{xsolution}

\paragraph{练习题 3}
\begin{xsolution}
\begin{enumerate}[label=(\arabic*)]
 \item 存在。例如
 \[
  f(x)=\tan\pi\left(x-\frac12\right),\qquad 0<x<1,
 \]
 连续且值域为 $\RR$。

 \item 不存在。若 $f\in C[0,1]$，由有界性定理 $f([0,1])$ 有界，不可能等于 $\RR$。

 \item 存在。例如
 \[
  f(x)=\frac{\sin(1/x)}x,\qquad 0<x\le1.
 \]
 它在 $(0,1]$ 连续。取使 $\sin(1/x)=1$ 与 $-1$ 的两列点趋于 $0$，可见函数在 $0$ 附近既无上界又无下界。连续像是区间，而该值域同时向上、向下无界，因此只能是整个 $\RR$。
\end{enumerate}
\end{xsolution}

\paragraph{练习题 4}
\begin{xsolution}
不一定。以 $[0,1]$ 为例，定义
\[
 f(x)=\begin{cases}
  x,&0<x<1,\\
  1,&x=0,\\
  0,&x=1.
 \end{cases}
\]
则 $f([0,1])=[0,1]$，确为闭区间，但 $f$ 在 $0,1$ 都不连续。因此“值域为闭区间”不能推出连续性。
\end{xsolution}

\paragraph{练习题 5}
\begin{xsolution}
设 $\lim_{x\to+\infty}f(x)=L$。若 $f\equiv L$，则最大值、最小值都能取到。以下设 $f$ 非常值。

若存在 $x_0$ 使 $f(x_0)>L$，取
\[
 \varepsilon=\frac{f(x_0)-L}{2}>0.
\]
可取 $M>x_0$，使 $x>M$ 时
\[
 f(x)<L+\varepsilon=\frac{L+f(x_0)}2<f(x_0).
\]
因此全局最大值只能出现在闭区间 $[a,M]$ 上，而连续函数在该闭区间上取到最大值，所以 $f$ 在 $[a,+\infty)$ 上取到最大值。

若不存在这样的 $x_0$，则 $f(x)\le L$ 对所有 $x$。因为 $f$ 非常值，存在 $x_0$ 使 $f(x_0)<L$。取
\[
 \varepsilon=\frac{L-f(x_0)}2>0.
\]
可取 $M>x_0$，使 $x>M$ 时
\[
 f(x)>L-\varepsilon=\frac{L+f(x_0)}2>f(x_0).
\]
因此全局最小值只能出现在闭区间 $[a,M]$ 上，并由最值定理在该区间取得。故至少最大值或最小值之一能够取到。
\end{xsolution}

\paragraph{练习题 6}
\begin{xsolution}
取中点 $x_0=(a+b)/2$。由
\[
 f(a^+)=f(b^-)=+\infty,
\]
可取
\[
 0<\delta<\frac{b-a}{2}
\]
使在 $(a,a+\delta)$ 与 $(b-\delta,b)$ 上都有
\[
 f(x)>f(x_0)+1.
\]
在闭区间 $[a+\delta,b-\delta]$ 上，$f$ 由最值定理取到最小值 $m$。因 $x_0$ 属于该闭区间，故 $m\le f(x_0)$；而两端邻近部分的函数值都大于 $f(x_0)+1>m$。因此 $m$ 也是整个 $(a,b)$ 上的最小值。
\end{xsolution}

\paragraph{练习题 7}
\begin{xsolution}
记共同的端点极限为 $L$。

若 $L\in\RR$，定义
\[
 F(a)=F(b)=L,\qquad F(x)=f(x)\quad(a<x<b).
\]
则 $F\in C[a,b]$，故取到最大值和最小值。若 $f\equiv L$，结论当然成立；否则存在一点函数值大于 $L$ 或小于 $L$。前一种情况下 $F$ 的最大值必在 $(a,b)$ 取得，后一种情况下最小值必在 $(a,b)$ 取得。

若 $L=+\infty$，由上一题 $f$ 取到最小值；若 $L=-\infty$，对 $-f$ 应用上一题，得到 $f$ 取到最大值。故各种情况下至少能取到最大值或最小值中的一个。
\end{xsolution}

\paragraph{练习题 8}
\begin{xsolution}
设 $m=f(a)$ 是 $f$ 的全局最小值，且 $m<a$。因为 $f(x)\to+\infty$ 当 $x\to\pm\infty$，连续函数的值域为
\[
 f(\RR)=[m,+\infty).
\]
特别地 $a\in[m,+\infty)=f(\RR)$。对复合函数 $F=f\circ f$，其取值为 $f$ 在 $[m,+\infty)$ 上的取值。由于 $a\in[m,+\infty)$ 且 $f(a)=m$ 是 $f$ 的全局最小值，故
\[
 \min_{x\in\RR}F(x)=m,
\]
而 $F(x)=m$ 当且仅当 $f(x)$ 是 $f$ 的某个最小值点；特别地，只要 $f(x)=a$ 就有 $F(x)=m$。

由 $f(a)=m<a$，再利用 $f(x)\to+\infty$ 当 $x\to-\infty$ 与 $x\to+\infty$，分别在 $a$ 的左、右两侧用介值定理，可得到
\[
 x_1<a<x_2,
 \qquad f(x_1)=f(x_2)=a.
\]
于是
\[
 F(x_1)=F(x_2)=f(a)=m,
\]
所以复合函数至少在两个不同点取到它的最小值。
\end{xsolution}

\paragraph{练习题 9}
\begin{xsolution}
对 Dirichlet 函数 $D$，有 $0\le D\le1$。每个有理点的函数值为 $1$，故每个有理点都是极大值点，同时也是最大值点；每个无理点函数值为 $0$，故每个无理点都是极小值点，同时也是最小值点。所有这些极值均非严格极值。

对 Riemann 函数 $R$，有 $R\ge0$。每个无理点函数值为 $0$，因此每个无理点都是极小值点和最小值点。若 $x_0=p/q$ 为既约有理数，$q>0$，则 $R(x_0)=1/q$。在 $x_0$ 的充分小邻域内，除 $x_0$ 外不存在分母不超过 $q$ 的既约有理数，所以邻域内其余点均有
\[
 R(x)<\frac1q=R(x_0).
\]
因此每个有理点都是严格极大值点。Riemann 函数的全局最大值为 $1$，恰在分母 $q=1$ 的有理点，即所有整数处取得；全局最小值为 $0$，在所有无理点取得。
\end{xsolution}

\subsection*{5.4 一致连续性与 Cantor 定理}

\subsubsection*{5.4.2 思考题}

\paragraph{思考题 1}
\begin{xsolution}
正面叙述为：存在 $\varepsilon_0>0$，使得对每个 $\delta>0$，总能找到 $x',x''\in I$ 满足
\[
 |x'-x''|<\delta,
 \qquad
 |f(x')-f(x'')|\ge\varepsilon_0.
\]
等价地，存在 $\varepsilon_0>0$ 与两列点 $x_n',x_n''\in I$，使
\[
 |x_n'-x_n''|\to0,
 \qquad
 |f(x_n')-f(x_n'')|\ge\varepsilon_0.
\]
\end{xsolution}

\paragraph{思考题 2}
\begin{xsolution}
错误。例如在 $[0,1)$ 上取
\[
 f(x)=\frac1{1-x}.
\]
它在 $[0,1)$ 上连续，但当 $x\to1^-$ 时无界。若它一致连续，则有界半开区间上的一致连续函数应有界，矛盾。也可直接取 $x_n=1-1/n$、$y_n=1-1/(2n)$ 构造反例。
\end{xsolution}

\paragraph{思考题 3}
\begin{xsolution}
错误。函数 $f(x)=1/x$ 在 $(0,1)$ 内每个闭子区间上都连续，但在 $(0,1)$ 上不一致连续。
\end{xsolution}

\paragraph{思考题 4}
\begin{xsolution}
正确。因为 $a<c<d<b$，$f$ 在闭区间 $[c,d]\subset(a,b)$ 上连续。由 Cantor 定理，$f$ 在 $[c,d]$ 上一致连续，所以其限制在 $(c,d)$ 上也一致连续。
\end{xsolution}

\subsubsection*{5.4.5 练习题}

\paragraph{练习题 1}
\begin{xsolution}
给定 $\varepsilon>0$。若 $L>0$，取
\[
 \delta=\frac\varepsilon L.
\]
当 $x_1,x_2\in I$ 且 $|x_1-x_2|<\delta$ 时，由 Lipschitz 条件
\[
 |f(x_1)-f(x_2)|\le L|x_1-x_2|<\varepsilon.
\]
所以 $f$ 在 $I$ 上一致连续。
\end{xsolution}

\paragraph{练习题 2}
\begin{xsolution}
这里 $a,b$ 为有限端点。取 $\varepsilon=1$。由一致连续性存在 $\delta>0$，使 $x,y\in(a,b)$ 且 $|x-y|<\delta$ 时
\[
 |f(x)-f(y)|<1.
\]
把有限区间 $(a,b)$ 分成有限个长度小于 $\delta$ 的小区间，并在每个非空小区间中选一点 $\xi_j$。若 $x$ 落在第 $j$ 个小区间，则
\[
 |f(x)|\le |f(\xi_j)|+1.
\]
有限个数 $|f(\xi_j)|+1$ 的最大值就是 $f$ 的一个全局界。因此 $f$ 有界。

若把端点允许为无穷，则结论一般不成立，例如 $f(x)=x$ 在 $\RR$ 上一致连续但无界。
\end{xsolution}

\paragraph{练习题 3}
\begin{xsolution}
\begin{enumerate}[label=(\arabic*)]
 \item 线性组合一定一致连续。事实上，给定 $\varepsilon>0$，分别利用 $f,g$ 的一致连续性使
 \[
  |a|\,|f(x)-f(y)|<\frac\varepsilon2,
  \qquad
  |b|\,|g(x)-g(y)|<\frac\varepsilon2
 \]
 （系数为零时相应项忽略），再取两个 $\delta$ 的较小者即可。

 乘积在一般区间上不一定一致连续。例如在 $I=\RR$ 上 $f(x)=g(x)=x$ 都一致连续，但 $fg=x^2$ 不一致连续。若另外知道 $I$ 有界，则 $f,g$ 因一致连续而有界，此时由
 \[
  |f(x)g(x)-f(y)g(y)|
  \le |f(x)|\,|g(x)-g(y)|+|g(y)|\,|f(x)-f(y)|
 \]
 可知乘积也一致连续。

 \item 一定一致连续。给定 $\varepsilon>0$，由 $g$ 在 $I_2$ 上一致连续，存在 $\eta>0$，使 $u,v\in I_2$ 且 $|u-v|<\eta$ 时
 \[
  |g(u)-g(v)|<\varepsilon.
 \]
 再由 $f$ 在 $I_1$ 上一致连续，存在 $\delta>0$，使 $x,y\in I_1$ 且 $|x-y|<\delta$ 时
 \[
  |f(x)-f(y)|<\eta.
 \]
 因 $f(I_1)\subset I_2$，故此时
 \[
  |g(f(x))-g(f(y))|<\varepsilon.
 \]
\end{enumerate}
\end{xsolution}

\paragraph{练习题 4}
\begin{xsolution}
设 $T>0$ 是 $f$ 的一个周期。由 Cantor 定理，$f$ 在闭区间 $[0,3T]$ 上一致连续。给定 $\varepsilon>0$，取相应的 $\delta>0$，并可令 $\delta<T$。

任取 $x,y\in\RR$ 且 $|x-y|<\delta$。可取整数 $k$，使
\[
 x-kT\in[T,2T).
\]
于是因为 $|x-y|<T$，有
\[
 y-kT\in(0,3T).
\]
利用周期性与 $[0,3T]$ 上的一致连续性，
\[
 |f(x)-f(y)|
 =|f(x-kT)-f(y-kT)|<\varepsilon.
\]
所以 $f$ 在整个 $\RR$ 上一致连续。
\end{xsolution}

\paragraph{练习题 5}
\begin{xsolution}
\begin{enumerate}[label=(\arabic*)]
 \item 写成
 \[
  f(x)=ax+b+r(x),
  \qquad r(x)=f(x)-(ax+b).
 \]
 则 $r\in C[0,+\infty)$ 且 $r(x)\to0$。由例题 5.4.6，$r$ 在 $[0,+\infty)$ 上一致连续；线性函数 $ax+b$ 满足 Lipschitz 条件，也一致连续。因此二者之和 $f$ 一致连续。

 \item 一般不成立。取 $f(x)=x^2$，则
 \[
  f(x)-(x^2)=0,
 \]
 完全满足相应的渐近条件，但 $x^2$ 在 $[0,+\infty)$ 上不一致连续。

 \item 若 $g$ 在 $[0,+\infty)$ 上一致连续，则 (1) 的结论仍成立。事实上，令
 \[
  r(x)=f(x)-g(x).
 \]
 因 $g$ 一致连续，所以 $g$ 连续，从而 $r\in C[0,+\infty)$；又由题设 $r(x)\to0$。例题 5.4.6 表明 $r$ 一致连续；再由 $g$ 一致连续，得到 $f=g+r$ 一致连续。

 在假定 $g$ 连续的条件下，这一条件也是必要的：若 $f$ 一致连续，则 $r=f-g$ 由上述理由一致连续，从而 $g=f-r$ 也一致连续。因此，此时结论成立的充分必要条件就是 $g$ 一致连续。
\end{enumerate}
\end{xsolution}

\paragraph{练习题 6}
\begin{xsolution}
这里 $n>1$ 不必是整数。取
\[
 x_m=m,
 \qquad
 y_m=(m^n+1)^{1/n}.
\]
则
\[
 y_m^n-x_m^n=1.
\]
又因 $0<1/n<1$，对 $t\ge0$ 有 $(1+t)^{1/n}\le1+t$，故
\[
 0<y_m-x_m
 =m\left[\left(1+\frac1{m^n}\right)^{1/n}-1\right]
 \le \frac1{m^{n-1}}\to0.
\]
因此两点距离趋于 $0$，而函数值之差恒为 $1$，与一致连续性的序列判据矛盾。所以 $x^n$ 在 $[0,+\infty)$ 上不一致连续。
\end{xsolution}

\paragraph{练习题 7}
\begin{xsolution}
在 $(1,+\infty)$ 上，设 $1<x<y$。由基本不等式 $\ln(1+t)\le t\ (t>-1)$，
\[
 0<\ln y-\ln x
 =\ln\left(1+\frac{y-x}{x}\right)
 \le\frac{y-x}{x}<y-x.
\]
故 $\ln x$ 在 $(1,+\infty)$ 上满足 Lipschitz 常数 $1$ 的估计，因而一致连续。

在 $(0,1)$ 上取
\[
 x_n=\frac1n,
 \qquad y_n=\frac1{2n}.
\]
则 $|x_n-y_n|\to0$，但
\[
 |\ln x_n-\ln y_n|=\ln2.
\]
故不一致连续。
\end{xsolution}

\paragraph{练习题 8}
\begin{xsolution}
\begin{enumerate}[label=(\arabic*)]
 \item 在 $(0,1)$ 上
 \[
  f(x)=\frac{|\sin x|}{x}=\frac{\sin x}{x}\longrightarrow1
  \qquad(x\to0^+),
 \]
 因而可连续延拓到 $[0,1]$；在 $(-1,0)$ 上
 \[
  f(x)=\frac{|\sin x|}{x}=-\frac{\sin x}{x}\longrightarrow-1
  \qquad(x\to0^-),
 \]
 也可连续延拓到 $[-1,0]$。由 Cantor 定理，两侧分别一致连续。

 但取 $x_n=-1/n,y_n=1/n$，则 $|x_n-y_n|\to0$，而
 \[
  f(x_n)\to-1,
  \qquad f(y_n)\to1,
 \]
 故在并集 $(-1,0)\cup(0,1)$ 上不一致连续。

 \item 不一定。令
 \[
  f(x)=\begin{cases}0,&a<x<b,\\1,&b\le x<c.
  \end{cases}
 \]
 则 $f$ 在 $(a,b)$ 与 $[b,c)$ 上分别为常值函数，故分别一致连续，但在 $b$ 处不连续，所以在 $(a,c)$ 上不可能一致连续。这里关键在于第一段不包含公共点 $b$；这与例题 5.4.4 的 $(a,b]$、$[b,c)$ 情形不同。
\end{enumerate}
\end{xsolution}

\paragraph{练习题 9}
\begin{xsolution}
\begin{enumerate}[label=(\arabic*)]
 \item $x\sin(1/x)$ 在 $x\to0^+$ 时趋于 $0$，故令 $f(0)=0$ 后可连续延拓到 $[0,1]$。由 Cantor 定理，它在 $(0,1)$ 上一致连续。

 \item 令 $f(0)=1$，则 $\sin x/x$ 连续延拓到 $[0,+\infty)$，而且 $f(x)\to0$ 当 $x\to+\infty$。由例题 5.4.6，它在 $[0,+\infty)$ 上一致连续，故在 $(0,+\infty)$ 上也一致连续。

 \item 不一致连续。取
 \[
  x_n=n\pi,
  \qquad y_n=n\pi+\frac1n.
 \]
 则 $|x_n-y_n|=1/n\to0$，而
 \[
  x_n\sin x_n=0,
 \]
 且
 \[
  |y_n\sin y_n|
  =\left(n\pi+\frac1n\right)\sin\frac1n
  \longrightarrow\pi.
 \]
 因此函数值之差不趋于 $0$，故 $x\sin x$ 在 $[0,+\infty)$ 上不一致连续。
\end{enumerate}
\end{xsolution}

\subsection*{5.5 单调函数}

\subsubsection*{5.5.2 练习题}

\paragraph{练习题 1}
\begin{xsolution}
任取 $x<y$，其中 $x,y\in(a,b)$。对每个 $t\in[x,y]$，题设给出一个邻域 $O(t)$，使 $f$ 在其中单调增加。这些邻域覆盖紧区间 $[x,y]$。由 Lebesgue 数性质，存在 $\delta>0$，使 $[x,y]$ 中距离小于 $\delta$ 的任意两点都落在某一个这样的邻域内。

取分点
\[
 x=t_0<t_1<\cdots<t_N=y,
 \qquad t_i-t_{i-1}<\delta.
\]
每对相邻分点同属某个 $O(t)$，而 $f$ 在该邻域单调增加，故
\[
 f(t_{i-1})\le f(t_i),\qquad i=1,\ldots,N.
\]
连锁得到 $f(x)\le f(y)$。由于 $x<y$ 任意，$f$ 在 $(a,b)$ 上单调增加。
\end{xsolution}

\paragraph{练习题 2}
\begin{xsolution}
任取 $x<y$。可取有理数列 $r_n\to x$、$s_n\to y$，并使对充分大的 $n$ 有
\[
 r_n<s_n,
 \qquad r_n,s_n\in[a,b].
\]
由题设
\[
 f(r_n)\le f(s_n).
\]
令 $n\to\infty$，利用 $f$ 在 $x,y$ 的连续性，得到
\[
 f(x)\le f(y).
\]
故 $f$ 在 $[a,b]$ 上单调增加。
\end{xsolution}

\paragraph{练习题 3}
\begin{xsolution}
若 $f$ 单调减少，则 $-f$ 单调增加，而且 $(-f)(x^+)=-g(x)$；因此只要证明单调增加的情形，便可对 $-f$ 应用同一结论得到 $g$ 右连续。以下设 $f$ 单调增加。此时
\[
 g(x)=f(x^+)=\inf_{t>x}f(t),
\]
且 $g$ 也单调增加。固定 $x_0$。给定 $\varepsilon>0$，由 $f(x_0^+)$ 的定义，可取 $y>x_0$ 使
\[
 f(y)<g(x_0)+\varepsilon.
\]
若 $x_0<x<y$，由单调性
\[
 g(x_0)\le g(x)\le f(y)<g(x_0)+\varepsilon.
\]
于是
\[
 \lim_{x\to x_0^+}g(x)=g(x_0),
\]
即 $g$ 在每一点右连续。
\end{xsolution}

\paragraph{练习题 4}
\begin{xsolution}
由加法方程先得
\[
 f(0)=0,
 \qquad f(-x)=-f(x),
 \qquad f(nx)=nf(x)\quad(n\in\ZZ).
\]
若 $f$ 单调增加，则对 $0<x<1/n$ 有
\[
 0=f(0)\le f(x)\le f(1/n)=\frac{f(1)}n\to0.
\]
若 $f$ 单调减少，同样得到
\[
 |f(x)|\le\frac{|f(1)|}{n}\qquad(0<x<1/n).
\]
所以无论哪一种单调性都有 $f(x)\to0$ 当 $x\to0$，即 $f$ 在 $0$ 连续。由
\[
 f(x+h)-f(x)=f(h)
\]
可知 $f$ 处处连续。连续加法函数满足
\[
 f(x)=f(1)x,
\]
故结论成立。
\end{xsolution}


\subsection*{5.7.2 参考题}

\subsubsection*{第一组参考题}

\paragraph{第一组参考题 1}
\begin{xsolution}
固定 $a\in(0,1)$，在 $[0,1-a]$ 上定义
\[
 g(x)=f(x+a)-f(x).
\]
则 $g$ 连续。由 $f\ge0$ 及 $f(0)=f(1)=0$，
\[
 g(0)=f(a)\ge0,
 \qquad
 g(1-a)=f(1)-f(1-a)=-f(1-a)\le0.
\]
若其中一个等于 $0$ 已得结论；否则由零点存在定理存在 $x_0\in(0,1-a)$ 使 $g(x_0)=0$，即
\[
 f(x_0)=f(x_0+a).
\]

去掉非负条件后结论一般不成立。例如取
\[
 f(x)=\sin2\pi x,
 \qquad a=\frac34.
\]
则 $f(0)=f(1)=0$。对 $x\in[0,1/4]$，令 $t=2\pi x\in[0,\pi/2]$，有
\[
 f\left(x+\frac34\right)-f(x)
 =-\cos t-\sin t<0.
\]
因此不存在 $x\in[0,1/4]$ 使 $f(x)=f(x+3/4)$。
\end{xsolution}

\paragraph{第一组参考题 2}
\begin{xsolution}
固定 $n\in\NN_+$。若 $n=1$，取 $\xi=0$ 即可。以下设 $n\ge2$，在 $[0,1-1/n]$ 上定义
\[
 g(x)=f\left(x+\frac1n\right)-f(x).
\]
若 $g$ 有零点即得结论。若无零点，由连续性知 $g$ 在整个区间恒正或恒负。但
\[
 \sum_{k=0}^{n-1}g\left(\frac{k}{n}\right)
 =\sum_{k=0}^{n-1}
 \left[f\left(\frac{k+1}{n}\right)-f\left(\frac{k}{n}\right)\right]
 =f(1)-f(0)=0,
\]
这与所有项同号矛盾。因此 $g$ 必有零点 $\xi$。
\end{xsolution}

\paragraph{第一组参考题 3}
\begin{xsolution}
令
\[
 F(x)=f(x)-x.
\]
由 $f([0,1])\subset[0,1]$，
\[
 F(0)=f(0)\ge0,
 \qquad F(1)=f(1)-1\le0.
\]
故存在 $x_1\in[0,1]$ 使 $F(x_1)=0$，即图像与 $y=x$ 相交。

再令
\[
 G(x)=f(x)+x-1.
\]
则
\[
 G(0)=f(0)-1\le0,
 \qquad G(1)=f(1)\ge0.
\]
故存在 $x_2\in[0,1]$ 使 $G(x_2)=0$，即 $f(x_2)=1-x_2$，图像与 $y=1-x$ 相交。
\end{xsolution}

\paragraph{第一组参考题 4}
\begin{xsolution}
这里把“极限存在”理解为存在有限实数极限。结论为：
\begin{enumerate}[label=(\arabic*)]
 \item $\lim_{x\to+\infty}f(x)$ 存在的充分必要条件是 $f$ 在某个尾区间 $[A,+\infty)$ 上有界；
 \item $\lim_{x\to0^+}f(x)$ 存在的充分必要条件是 $f$ 在某个 $(0,A]$ 上有界。
\end{enumerate}
必要性都由有限极限的局部有界性得到。先证第一条的充分性。

设 $f$ 在 $[A,+\infty)$ 上有界。若 $f(x)$ 在 $+\infty$ 处没有有限极限，则其尾部下极限与上极限满足
\[
 \alpha=\varliminf_{x\to+\infty}f(x)
 <
 \beta=\varlimsup_{x\to+\infty}f(x).
\]
取 $c\in(\alpha,\beta)$。于是无论把起点取得多大，后面总能找到函数值小于 $c$ 的点和函数值大于 $c$ 的点。可归纳取严格递增的点列
\[
 u_1<v_1<u_2<v_2<\cdots\to+\infty
\]
使 $f(u_j)<c<f(v_j)$（必要时交换两类点的次序）。由连续性，每个 $[u_j,v_j]$ 中都有一点 $z_j$ 满足 $f(z_j)=c$。这使方程 $f(x)=c$ 有无限多个实根，与题设矛盾。因此 $\alpha=\beta$，有限极限存在。

再证第二条的充分性。若 $f$ 在 $(0,A]$ 上有界，令
\[
 h(t)=f(1/t),
 \qquad t\in[1/A,+\infty).
\]
则 $h$ 连续且有界；并且对每个 $c\in\RR$，方程 $h(t)=c$ 的根数不超过方程 $f(x)=c$ 的根数，因而也是有限的。由刚证得的第一条，$\lim_{t\to+\infty}h(t)$ 存在。令 $x=1/t$，即得 $\lim_{x\to0^+}f(x)$ 存在。
\end{xsolution}

\paragraph{第一组参考题 5}
\begin{xsolution}
\begin{enumerate}[label=(\arabic*)]
 \item 对 $x<y$，由题设
 \[
  f(y)-f(x)\le |f(y)-f(x)|\le k(y-x),
 \]
 故
 \[
  [ky-f(y)]-[kx-f(x)]
  =k(y-x)-[f(y)-f(x)]\ge0.
 \]
 因此 $kx-f(x)$ 单调增加。

 \item 令 $F(x)=f(x)-x$。对 $x<y$，
 \[
  F(y)-F(x)=f(y)-f(x)-(y-x)
  \le -(1-k)(y-x)<0,
 \]
 所以 $F$ 严格减少，零点至多一个。

 又由 Lipschitz 条件
 \[
  f(x)\le f(0)+kx\quad(x>0),
 \]
 从而
 \[
  F(x)\le f(0)-(1-k)x\to-\infty
  \qquad(x\to+\infty).
 \]
 对 $x<0$，
 \[
  f(x)\ge f(0)-k|x|=f(0)+kx,
 \]
 故
 \[
  F(x)\ge f(0)+(k-1)x\to+\infty
  \qquad(x\to-\infty).
 \]
 $F$ 连续，因此由零点存在定理有且仅有一个 $\xi$ 满足 $F(\xi)=0$，即 $f(\xi)=\xi$。
\end{enumerate}
\end{xsolution}

\paragraph{第一组参考题 6}
\begin{xsolution}
\begin{enumerate}[label=(\arabic*)]
 \item 对 $n>1$，函数 $f_n(x)=x^n+x$ 在 $(0,+\infty)$ 上严格增加。并且
 \[
  f_n(1/2)=2^{-n}+\frac12<1,
  \qquad f_n(1)=2>1.
 \]
 由零点存在定理，$f_n(x)=1$ 在 $(1/2,1)$ 内有根；严格增加性又保证根唯一。

 \item 记唯一根为 $c_n$。因 $0<c_n<1$，有
 \[
  f_{n+1}(c_n)=c_n^{n+1}+c_n<c_n^n+c_n=1.
 \]
 而 $f_{n+1}$ 严格增加，所以
 \[
  c_{n+1}>c_n.
 \]
 因此 $\{c_n\}$ 单调增加且以 $1$ 为上界，故存在极限 $L\le1$。若 $L<1$，取 $q$ 使 $L<q<1$，则充分大时 $c_n<q$，从而 $c_n^n\le q^n\to0$。由方程
 \[
  c_n^n=1-c_n
 \]
 令 $n\to\infty$ 得 $0=1-L$，矛盾。因此
 \[
  \boxed{\lim_{n\to\infty}c_n=1}.
 \]
\end{enumerate}
\end{xsolution}

\paragraph{第一组参考题 7}
\begin{xsolution}
固定 $a\in[0,1]$ 与 $\varepsilon>0$。取 $N$ 使 $1/N<\varepsilon$。有限并
\[
 E=A_1\cup\cdots\cup A_N
\]
仍为有限集。可取 $\delta>0$，使去心邻域
\[
 0<|x-a|<\delta
\]
中没有 $E$ 的点。于是这样的 $x$ 若属于某个 $A_n$，必有 $n>N$，从而
\[
 |f(x)|=\frac1n<\frac1N<\varepsilon;
\]
若不属于任何 $A_n$，则 $f(x)=0$。所以
\[
 \boxed{\lim_{x\to a}f(x)=0}
\]
对每个 $a\in[0,1]$ 都成立。
\end{xsolution}

\paragraph{第一组参考题 8}
\begin{xsolution}
记
\[
 g(x)=\lim_{t\to x}f(t).
\]
固定 $x_0\in I$。给定 $\varepsilon>0$，由 $f(t)\to g(x_0)$ 当 $t\to x_0$，可取 $\delta>0$，使
\[
 0<|t-x_0|<\delta
 \quad\Longrightarrow\quad
 |f(t)-g(x_0)|<\frac\varepsilon2.
\]
若 $|x-x_0|<\delta/2$，则取 $t\to x$ 且 $t\ne x$ 时，可令 $t$ 同时落在 $O_\delta(x_0)$ 内。于是
\[
 |f(t)-g(x_0)|<\frac\varepsilon2.
\]
令 $t\to x$，得到
\[
 |g(x)-g(x_0)|\le\frac\varepsilon2<\varepsilon.
\]
所以 $g$ 在 $x_0$ 连续。$x_0$ 任意，故 $g\in C(I)$。
\end{xsolution}

\paragraph{第一组参考题 9}
\begin{xsolution}
若 $\lambda=0$，任取 $x_n\to+\infty$ 即可。先设 $\lambda>0$。对于 $\lambda<0$，令 $\mu=-\lambda>0$；若已对 $\mu$ 找到 $y_n\to+\infty$ 使 $f(y_n+\mu)-f(y_n)\to0$，再取 $x_n=y_n+\mu$，便有
\[
 f(x_n+\lambda)-f(x_n)=f(y_n)-f(y_n+\mu)\to0.
\]
因此只需证明 $\lambda>0$ 的情形。

令
\[
 g(x)=f(x+\lambda)-f(x).
\]
若结论不成立，则存在 $\varepsilon_0>0$ 和 $M$，使对所有 $x>M$ 都有
\[
 |g(x)|\ge\varepsilon_0.
\]
因为 $g$ 连续且在连通区间 $(M,+\infty)$ 上不取 $0$，它在该区间恒正或恒负。若 $g(x)\ge\varepsilon_0$，则对固定 $x>M$，
\[
 f(x+n\lambda)-f(x)
 =\sum_{j=0}^{n-1}g(x+j\lambda)
 \ge n\varepsilon_0,
\]
与 $f$ 有界矛盾；若 $g\le-\varepsilon_0$ 同理矛盾。因此对每个 $n$ 都能取 $x_n>n$ 使
\[
 |f(x_n+\lambda)-f(x_n)|<\frac1n.
\]
于是 $x_n\to+\infty$ 且所要求的差趋于 $0$。
\end{xsolution}

\paragraph{第一组参考题 10}
\begin{xsolution}
任取 $x<y$，把 $[x,y]$ 等分为 $n$ 段，分点记为
\[
 x=x_0<x_1<\cdots<x_n=y,
 \qquad x_j-x_{j-1}=\frac{y-x}{n}.
\]
由题设及三角不等式，
\[
 \begin{aligned}
 |f(y)-f(x)|
 &\le\sum_{j=1}^n|f(x_j)-f(x_{j-1})|\\
 &\le nM\left(\frac{y-x}{n}\right)^\alpha
 =M(y-x)^\alpha n^{1-\alpha}.
 \end{aligned}
\]
若 $\alpha>1$，令 $n\to\infty$，右端趋于 $0$，故 $f(y)=f(x)$。任意两点函数值相同，所以 $f$ 为常值函数。
\end{xsolution}

\paragraph{第一组参考题 11}
\begin{xsolution}
先定义
\[
 h(x)=\begin{cases}
 x,&x\in\QQ,\\
 1-x,&x\notin\QQ,
 \end{cases}
 \qquad 0\le x\le1.
\]
$h$ 的值域恰为 $[0,1]$，并且除 $x=1/2$ 外处处不连续；在 $1/2$ 处两类趋近的函数值都趋于 $1/2$。

现只交换 $h$ 在 $0$ 与 $1/2$ 的两个函数值，定义
\[
 f(0)=\frac12,
 \qquad f\left(\frac12\right)=0,
\]
其余点令 $f(x)=h(x)$。交换两个值不改变值域，所以 $f([0,1])=[0,1]$。对 $x_0\ne1/2$，有理、无理数列趋于 $x_0$ 时函数值分别趋于 $x_0$ 与 $1-x_0$，两者不同，故不连续；在 $x_0=1/2$，极限为 $1/2$ 而函数值为 $0$，也不连续。因此 $f$ 处处不连续而值域为区间。
\end{xsolution}

\paragraph{第一组参考题 12}
\begin{xsolution}
由 Cantor 定理，$f$ 在 $[a,b]$ 上一致连续。给定 $\varepsilon>0$，取 $\delta>0$，使
\[
 |x-y|<\delta
 \quad\Longrightarrow\quad
 |f(x)-f(y)|<\varepsilon.
\]
取分点
\[
 a=x_0<x_1<\cdots<x_N=b
\]
使每段长度小于 $\delta$，并令 $L$ 在每个 $[x_{j-1},x_j]$ 上为连接
\[
 (x_{j-1},f(x_{j-1})),\qquad (x_j,f(x_j))
\]
的线性函数。若 $x\in[x_{j-1},x_j]$，则存在 $0\le\theta\le1$ 使
\[
 L(x)=(1-\theta)f(x_{j-1})+\theta f(x_j).
\]
于是
\[
 \begin{aligned}
 |f(x)-L(x)|
 &\le(1-\theta)|f(x)-f(x_{j-1})|
   +\theta|f(x)-f(x_j)|\\
 &<\varepsilon.
 \end{aligned}
\]
故得到所要求的分段线性函数。
\end{xsolution}

\paragraph{第一组参考题 13}
\begin{xsolution}
这里无符号的 $\infty$ 表示无穷大，即题设为
\[
 |f(f(x))|\to+\infty,
\]
所需证明的结论为 $|f(x)|\to+\infty$。

反设结论不成立，则存在常数 $M>0$ 和数列 $x_n\to+\infty$，使
\[
 |f(x_n)|\le M,
 \qquad n\in\NN_+.
\]
由于 $f$ 在闭区间 $[-M,M]$ 上连续，存在 $K>0$ 使
\[
 |f(y)|\le K,
 \qquad |y|\le M.
\]
于是
\[
 |f(f(x_n))|\le K,
\]
这与 $|f(f(x))|\to+\infty$ 矛盾。因此
\[
 |f(x)|\to+\infty.
\]
\end{xsolution}

\paragraph{第一组参考题 14}
\begin{xsolution}
题设立即给出一一性：若 $f(x_1)=f(x_2)$，则
\[
 0\ge k|x_1-x_2|,
\]
故 $x_1=x_2$。连续的一一函数在区间 $\RR$ 上严格单调。

又取 $x_2=0$ 得
\[
 |f(x)-f(0)|\ge k|x|,
\]
从而
\[
 |f(x)|\ge k|x|-|f(0)|\to+\infty
 \qquad(|x|\to\infty).
\]
若 $f$ 严格增加，则不可能在 $+\infty$ 方向趋于 $-\infty$，也不可能在 $-\infty$ 方向趋于 $+\infty$，故
\[
 \lim_{x\to-\infty}f(x)=-\infty,
 \qquad
 \lim_{x\to+\infty}f(x)=+\infty.
\]
若严格减少则两式符号相反。两种情况下由介值定理都得到
\[
 f(\RR)=\RR.
\]
\end{xsolution}

\paragraph{第一组参考题 15}
\begin{xsolution}
设 $P$ 为 $f$ 的所有正周期组成的集合。先证
\[
 \tau=\inf P>0.
\]
若不然，可取周期 $T_n\downarrow0$。任取 $x,y\in\RR$，对每个 $n$ 取整数 $m_n$，使
\[
 |x+m_nT_n-y|\le\frac{T_n}{2}.
\]
由周期性 $f(x+m_nT_n)=f(x)$，而 $x+m_nT_n\to y$。连续性给出
\[
 f(x)=\lim_{n\to\infty}f(x+m_nT_n)=f(y).
\]
这会使 $f$ 成为常值函数，与题设矛盾。因此 $\tau>0$。

按下确界定义可取 $T_n\in P$ 使 $T_n\to\tau$。对任意 $x$，
\[
 f(x+\tau)
 =\lim_{n\to\infty}f(x+T_n)
 =f(x),
\]
故 $\tau$ 本身也是周期。由定义它就是最小正周期。

若去掉连续性，结论不成立。Dirichlet 函数 $D$ 的每个正有理数都是周期，因为平移有理数不改变一个数的有理性；但正有理数没有最小元，所以 $D$ 没有最小正周期。
\end{xsolution}

\paragraph{第一组参考题 16}
\begin{xsolution}
设 $f$ 的 Lipschitz 常数为 $L$。对 $x\ge a$，
\[
 |f(x)|\le |f(a)|+L(x-a)\le Lx+B,
 \qquad B=|f(a)|+La.
\]
令 $g(x)=f(x)/x$。对 $x,y\ge a$，
\[
 \begin{aligned}
 |g(x)-g(y)|
 &\le\frac{|f(x)-f(y)|}{x}
 +|f(y)|\left|\frac1x-\frac1y\right|\\
 &\le\frac{L}{a}|x-y|
 +(Ly+B)\frac{|x-y|}{xy}\\
 &\le\left[\frac{L}{a}+\frac{L+B/a}{a}\right]|x-y|.
 \end{aligned}
\]
所以 $g$ 本身满足 Lipschitz 条件，从而在 $[a,+\infty)$ 上一致连续。
\end{xsolution}

\paragraph{第一组参考题 17}
\begin{xsolution}
若 $f$ 一致连续，而 $x_n-y_n\to0$，则给定 $\varepsilon>0$，取一致连续性中的 $\delta$。充分大时 $|x_n-y_n|<\delta$，于是
\[
 |f(x_n)-f(y_n)|<\varepsilon,
\]
故差趋于 $0$。

反之，若 $f$ 不一致连续，则存在 $\varepsilon_0>0$，对每个 $n$ 可取 $x_n,y_n\in I$ 使
\[
 |x_n-y_n|<\frac1n,
 \qquad
 |f(x_n)-f(y_n)|\ge\varepsilon_0.
\]
于是 $x_n-y_n\to0$，但函数值之差不趋于 $0$，与题设的序列条件矛盾。因此该条件也充分。
\end{xsolution}

\paragraph{第一组参考题 18}
\begin{xsolution}
若 $f$ 一致连续，设 $\{x_n\}\subset I$ 为基本数列。给定 $\varepsilon>0$，取一致连续性中的 $\delta$；再由 $\{x_n\}$ 为基本数列，取 $N$ 使 $m,n>N$ 时 $|x_m-x_n|<\delta$。于是
\[
 |f(x_m)-f(x_n)|<\varepsilon,
\]
故 $\{f(x_n)\}$ 也是基本数列。

反之，若 $f$ 不一致连续，则存在 $\varepsilon_0>0$ 以及 $x_n,y_n\in I$，使
\[
 |x_n-y_n|<\frac1n,
 \qquad |f(x_n)-f(y_n)|\ge\varepsilon_0.
\]
因 $I$ 有界，$\{x_n\}$ 有收敛子列 $x_{n_k}\to\xi$；相应地 $y_{n_k}\to\xi$。交错排列
\[
 z_{2k-1}=x_{n_k},
 \qquad z_{2k}=y_{n_k},
 \qquad k\in\NN_+,
\]
得到一个 $I$ 中的基本数列。但
\[
 |f(z_{2k-1})-f(z_{2k})|
 \ge\varepsilon_0,
\]
所以 $\{f(z_n)\}$ 不是基本数列，矛盾。故 $f$ 必一致连续。
\end{xsolution}

\paragraph{第一组参考题 19}
\begin{xsolution}
给定 $\varepsilon>0$。由 $f$ 在 $[0,+\infty)$ 上一致连续，存在 $\delta>0$，使
\[
 |u-v|<\delta
 \Longrightarrow
 |f(u)-f(v)|<\frac\varepsilon2.
\]
在 $[0,1]$ 中选有限个点 $t_1,\ldots,t_m$，使每个 $t\in[0,1]$ 与某个 $t_j$ 的距离小于 $\delta$。对每个 $j$，由
\[
 f(t_j+n)\to0
\]
可取同一个足够大的整数 $N$，使当 $n\ge N$ 时
\[
 |f(t_j+n)|<\frac\varepsilon2,
 \qquad j=1,\ldots,m.
\]
任取充分大的 $x$，写成 $x=n+t$，其中 $n\in\NN$、$t\in[0,1)$，并令 $n\ge N$。选 $t_j$ 使 $|t-t_j|<\delta$，则
\[
 |f(x)|
 \le |f(n+t)-f(n+t_j)|+|f(n+t_j)|
 <\varepsilon.
\]
所以 $f(x)\to0$ 当 $x\to+\infty$。
\end{xsolution}

\paragraph{第一组参考题 20}
\begin{xsolution}
由一致连续性，对 $\varepsilon=1$ 可取 $\delta>0$，使
\[
 |u-v|<\delta
 \Longrightarrow
 |f(u)-f(v)|<1.
\]
固定 $x\in\RR$。把从 $0$ 到 $x$ 的线段分成 $N$ 段，使每段长度小于 $\delta$，并可取
\[
 N\le\frac{2|x|}{\delta}+1.
\]
沿分点逐段使用三角不等式，得到
\[
 |f(x)-f(0)|<N
 \le\frac{2|x|}{\delta}+1.
\]
因此
\[
 |f(x)|
 \le \frac{2}{\delta}|x|+|f(0)|+1.
\]
取
\[
 a=\frac2\delta\ge0,
 \qquad b=|f(0)|+1\ge0,
\]
即得所需估计。
\end{xsolution}


\subsubsection*{第二组参考题}

\paragraph{第二组参考题 1}
\begin{xsolution}
下面都只使用题目允许的朴素几何观点，并把面积随直线位置、方向连续变化这一事实作为连续性工具。

\begin{enumerate}[label=(\arabic*)]
 \item 固定一个方向，让一条与该方向垂直的直线从煎饼一侧平行移动到另一侧。记直线某一侧的面积减去另一侧的面积为 $F(t)$。当直线在煎饼一侧之外时，$F$ 的符号为负；移到另一侧之外时，符号为正。$F(t)$ 连续，故由零点存在定理存在位置使 $F(t)=0$，即一刀把面积二等分。

 \item 对每个方向 $\theta$，先取一条方向为 $\theta$ 的直线 $L_\theta$，使第一个煎饼被二等分。若这样的直线不唯一，可在所有二等分位置组成的区间中取中点，从而使选择随 $\theta$ 连续。令
 \[
  H(\theta)=\text{第二个煎饼在 }L_\theta\text{ 两侧的面积差}.
 \]
 当方向反转为 $\theta+\pi$ 时，同一条无向直线的两侧互换，故
 \[
  H(\theta+\pi)=-H(\theta).
 \]
 $H$ 连续，因此在 $[\theta,\theta+\pi]$ 中必有 $H=0$ 的方向。相应直线同时二等分两个煎饼。

 \item 一般不能。取三个互不共线的圆盘。任何把一个圆盘面积二等分的直线必须经过圆心，因此若一条直线同时二等分三个圆盘，它必须同时经过三个不共线圆心，这是不可能的。

 \item 可以。对每个方向 $\theta$，分别取方向为 $\theta$ 和 $\theta+\pi/2$ 的两条面积二等分线 $L_\theta,M_\theta$；若某方向的二等分线不唯一，就取所有二等分位置组成区间的中点。这样两条线可随 $\theta$ 连续选取。它们把平面分成四个象限，按逆时针记煎饼落在四个象限中的面积为 $A_1,A_2,A_3,A_4$。因为 $L_\theta$、$M_\theta$ 各自都二等分总面积，
 \[
  A_1+A_2=A_3+A_4,
  \qquad
  A_2+A_3=A_4+A_1.
 \]
 因而 $A_1=A_3$、$A_2=A_4$。令
 \[
  G(\theta)=A_1-A_2.
 \]
 当 $\theta$ 增加 $\pi/2$ 时，两条互相垂直的二等分线交换角色，四个象限循环移动一格，因此
 \[
  G\left(\theta+\frac\pi2\right)=A_2-A_3=-G(\theta).
 \]
 面积对方向连续，所以 $G$ 连续。由零点存在定理存在 $\theta_0$ 使 $G(\theta_0)=0$。这时 $A_1=A_2=A_3=A_4$，而四者之和为总面积，故每块面积均为总面积的 $1/4$。

 \item 设运动员在时刻 $t$ 已跑过的路程为 $s(t)$。则 $s$ 连续，$s(0)=0,s(10)=100$。在 $[0,9]$ 上令
 \[
  g(t)=s(t+1)-s(t)-10.
 \]
 若不存在 $g(t)=0$，因 $g$ 连续，它必在整个 $[0,9]$ 上恒正或恒负。但
 \[
  \sum_{k=0}^{9}g(k)
  =s(10)-s(0)-100=0,
 \]
 与所有项同号矛盾。所以存在 $t_0\in[0,9]$ 使
 \[
  s(t_0+1)-s(t_0)=10,
 \]
 即有一段 $10\,\mathrm m$ 的路程恰用 $1\,\mathrm s$ 完成。

 \item 固定方台中心位置，只让方台绕竖直轴转动。把一条对角线两端两只脚的“着地高度和”减去另一条对角线两端两只脚的“着地高度和”，记为 $H(\theta)$。地面高度连续时，$H$ 随转角连续。转过 $\pi/2$ 后两条对角线交换，故
 \[
  H\left(\theta+\frac\pi2\right)=-H(\theta).
 \]
 因而某个角度有 $H=0$。此时四个脚点共面，刚性方台四脚可同时接触地面，于是不再摇晃。

 \item 对每个方向 $\theta$，作两条方向为 $\theta$ 的平行支撑线夹住闭曲线，记两线距离即该方向的宽度为 $w(\theta)$。光滑闭曲线的支撑线与曲线相切，且 $w(\theta)$ 连续。考虑
 \[
  H(\theta)=w(\theta)-w\left(\theta+\frac\pi2\right).
 \]
 则
 \[
  H\left(\theta+\frac\pi2\right)
  =-H(\theta).
 \]
 故存在 $\theta_0$ 使 $H(\theta_0)=0$。于是方向 $\theta_0$ 与垂直方向的两对支撑线距离相等，四条支撑线围成一个正方形，并且四边都与曲线相切。
\end{enumerate}
\end{xsolution}

\paragraph{第二组参考题 2}
\begin{xsolution}
把 $f$ 以周期 $n$ 连续延拓到 $\RR$，记延拓仍为 $f$。对整数 $k=1,\ldots,n-1$ 定义
\[
 h_k(x)=f(x+k)-f(x).
\]
令 $r=n/\gcd(n,k)$。由于在模 $n$ 意义下
\[
 x,x+k,\ldots,x+(r-1)k
\]
构成一个循环，有
\[
 \sum_{j=0}^{r-1}h_k(x+jk)=0.
\]
因此若 $h_k$ 不恒为零，它必取正值和负值，故在一个周期圆周上至少有两个零点；若恒为零，则零点当然更多。

现在只取
\[
 1\le k<\frac n2.
\]
对每个这样的 $k$，至少可得到两个不同的无序点对 $\{x,x+k\}$（模 $n$），其两点函数值相同。不同 $k$ 对应的圆周短距离不同，所以这些点对互不相同。若 $n$ 为偶数，再取 $k=n/2$。这时
\[
 h_{n/2}\left(x+\frac n2\right)=-h_{n/2}(x),
\]
故至少得到一个这样的点对。

于是当 $n$ 为奇数时得到
\[
 2\cdot\frac{n-1}{2}=n-1
\]
对；当 $n$ 为偶数时得到
\[
 2\left(\frac n2-1\right)+1=n-1
\]
对。把圆周上的各点取在 $[0,n)$ 的代表元中，每一对的两个代表元之差为非零整数。最后再加入端点对 $(0,n)$，由 $f(0)=f(n)$，总共至少得到 $n$ 对不同的 $(x,y)$。
\end{xsolution}

\paragraph{第二组参考题 3}
\begin{xsolution}
记
\[
 L(x)=\lim_{t\to x}f(t).
\]
由第一组参考题 8 的结论，$L\in C[a,b]$。

\begin{enumerate}[label=(\arabic*)]
 \item 固定 $\varepsilon>0$。若满足
 \[
  |L(x)-f(x)|>\varepsilon
 \]
 的点有无限多个，则可取互异点列 $x_n$，由闭区间的凝聚性再取子列（仍记为 $x_n$）使 $x_n\to x_0\in[a,b]$。

 因为 $f$ 在 $x_0$ 处有极限，且 $x_n\ne x_0$，所以
 \[
  f(x_n)\to L(x_0).
 \]
 又因 $L$ 连续，
 \[
  L(x_n)\to L(x_0).
 \]
 因此
 \[
  |L(x_n)-f(x_n)|\to0,
 \]
 与它始终大于 $\varepsilon$ 矛盾。故这样的点至多有限个。

 \item 点 $x$ 是可去间断点当且仅当 $f(x)\ne L(x)$。于是全部间断点集合为
 \[
  \bigcup_{m=1}^\infty
  \left\{x\in[a,b]:|L(x)-f(x)|>\frac1m\right\}.
 \]
 由 (1)，每个括号中的集合有限，故该并集至多可列。
\end{enumerate}
\end{xsolution}

\paragraph{第二组参考题 4}
\begin{xsolution}
若区间 $I$ 的每一点都是 $f$ 的可去间断点，记
\[
 L(x)=\lim_{t\to x}f(t).
\]
对任意有界闭子区间 $[c,d]\subset I$，由上一题，集合
\[
 E_m=\left\{x\in[c,d]:|L(x)-f(x)|>\frac1m\right\}
\]
都是有限集。而每一点都是可去间断点，故 $f(x)\ne L(x)$，从而
\[
 [c,d]=\bigcup_{m=1}^\infty E_m.
\]
右边是可列集，左边却不可列，矛盾。因此不存在这样的函数。
\end{xsolution}

\paragraph{第二组参考题 5}
\begin{xsolution}
\begin{enumerate}[label=(\arabic*)]
 \item 不存在。反设连续函数 $f\colon\RR\to\RR$ 对每个 $c$ 恰有两个原像。取某个 $c$，设两个原像为 $a<b$。在 $(a,b)$ 上 $f-c$ 不取零值，故恒正或恒负。

 若 $f>c$ 于 $(a,b)$，在 $[a,b]$ 上取最大值 $M>c$。取 $d\in(c,M)$。由介值定理，方程 $f(x)=d$ 在 $(a,b)$ 至少有两个解，因此按题设恰有这两个解。若在 $(-\infty,a)$ 或 $(b,+\infty)$ 中某处有 $f>d$，从该点到 $a$ 或 $b$ 再用介值定理会产生第三个 $d$ 的原像，矛盾。因此区间外有 $f<d$，而区间内 $f\le M$，所以 $f$ 在 $\RR$ 上有上界。这与每个任意大的实数都应成为某个函数值矛盾。若 $f<c$ 于 $(a,b)$，同理得到全局下界，也矛盾。

 \item 存在。定义连续分段线性函数 $F$，使对每个 $m\in\ZZ$，
 \[
  F(3m)=m,
  \quad F(3m+1)=m+1,
  \quad F(3m+2)=m,
  \quad F(3m+3)=m+1,
 \]
 并在相邻整数间线性插值。若 $c\in(m,m+1)$，水平线 $y=c$ 恰与区间
 \[
  (3m,3m+1),\quad(3m+1,3m+2),\quad(3m+2,3m+3)
 \]
 上的三条线段各交一次，所以恰有三个解。若 $c=m$ 是整数，则原像恰为
 \[
  3m-2,\quad3m,\quad3m+2,
 \]
 仍恰有三个。因此所求函数存在。
\end{enumerate}
\end{xsolution}

\paragraph{第二组参考题 6}
\begin{xsolution}
方程为
\[
 f(x+y^n)=f(x)+(f(y))^n,
 \qquad x,y\in\RR.
\]
令 $y=0$ 得 $f(0)=0$；再令 $x=0$ 得
\[
 f(y^n)=f(y)^n.
\]
于是
\[
 f(x+y^n)=f(x)+f(y^n).
\]

若 $n$ 为奇数，则 $y^n$ 遍历 $\RR$，故 $f$ 是加法函数。若 $n$ 为偶数，则上式先给出
\[
 f(x+t)=f(x)+f(t),\qquad t\ge0.
\]
取 $x=-t$ 得 $f(-t)=-f(t)$，于是同样可推广到任意实数增量，故 $f$ 仍是加法函数。

当 $n=1$ 时，原方程正是 Cauchy 加法方程，因此所有加法函数均为解。

以下设 $n>1$。利用恒等式
\[
 n!x_1\cdots x_n
 =\sum_{S\subset\{1,\ldots,n\}}
 (-1)^{n-|S|}
 \left(\sum_{i\in S}x_i\right)^n,
\]
对两边应用加法函数 $f$，并利用 $f(u^n)=f(u)^n$，得到
\[
 f(x_1\cdots x_n)=f(x_1)\cdots f(x_n).
\]
令 $c=f(1)$，则
\[
 c=c^n.
\]
若 $c=0$，在上式中令 $x_2=\cdots=x_n=1$ 即得 $f(x)=0$。

若 $c=1$，则 $f$ 是保持乘法的含幺实数域自同态。特别地，对 $u\ge0$，写 $u=v^2$，有
\[
 f(u)=f(v)^2\ge0.
\]
因此 $x<y$ 时 $f(y)-f(x)=f(y-x)\ge0$，故 $f$ 单调。单调加法函数连续，从而
\[
 f(x)=x.
\]

当 $n$ 为奇数且 $c=-1$ 时，令 $g=-f$。因为 $n$ 为奇数，$g$ 仍为加法且满足 $g(u^n)=g(u)^n$，并有 $g(1)=1$，由上一段 $g(x)=x$，故 $f(x)=-x$。

综上：
\[
\begin{cases}
 n=1:& f\text{ 为任意加法函数};\\
 n\ge2\text{ 为偶数}:& f\equiv0\text{ 或 }f(x)=x;\\
 n\ge3\text{ 为奇数}:& f\equiv0,\ f(x)=x\text{ 或 }f(x)=-x.
\end{cases}
\]
直接代入可验证这些函数均满足原方程。
\end{xsolution}

\paragraph{第二组参考题 7}
\begin{xsolution}
由
\[
 f(f(x))=x
\]
可知 $f$ 是一一映射。连续的一一函数在区间 $[0,1]$ 上必严格单调。又 $f(0)=0<f(1)=1$，所以 $f$ 严格增加。

若存在 $x$ 使 $f(x)>x$，由严格增加性，
\[
 f(f(x))>f(x),
\]
即 $x>f(x)$，矛盾。若 $f(x)<x$，则同理
\[
 f(f(x))<f(x),
\]
即 $x<f(x)$，仍矛盾。因此对每个 $x$ 只能有
\[
 f(x)=x.
\]
\end{xsolution}

\paragraph{第二组参考题 8}
\begin{xsolution}
若 $k=0$，取 $f\equiv0$ 即有连续解。

设 $k\ne0$ 且存在连续解。函数 $kx^9$ 为一一函数，所以若 $f(x_1)=f(x_2)$，再作用一次 $f$ 得
\[
 kx_1^9=kx_2^9,
\]
从而 $x_1=x_2$。故 $f$ 连续且一一，因而严格单调。不论 $f$ 严格增加还是严格减少，复合函数 $f\circ f$ 都严格增加。因此 $kx^9$ 必严格增加，只能有
\[
 k>0.
\]

反之，对任意 $k>0$，令
\[
 c=k^{1/4}>0,
 \qquad f(x)=cx^3.
\]
则
\[
 f(f(x))=c(cx^3)^3=c^4x^9=kx^9.
\]
故连续解存在。综上，充分必要条件为
\[
 \boxed{k\ge0}.
\]
\end{xsolution}

\paragraph{第二组参考题 9}
\begin{xsolution}
由
\[
 f(2x-f(x))=x
\]
知 $f$ 是满射；题设又说 $f$ 一一，所以 $f$ 是双射，并且
\[
 f^{-1}(x)=2x-f(x).
\]
若 $f$ 严格减少，则 $f^{-1}$ 也严格减少，而 $2x-f(x)$ 却严格增加，矛盾。因此连续一一函数 $f$ 必严格增加。

将上式中的 $x$ 换成 $f(t)$，得到
\[
 t=2f(t)-f(f(t)),
\]
即
\[
 f(f(t))-f(t)=f(t)-t.
\]
记
\[
 T(t)=f(t)-t,
\]
则
\[
 T(f(t))=T(t).
\]
因为 $f$ 为双射，这个等式可沿正、负方向迭代，所以若 $T(x)=c$，则对所有整数 $m$ 都有
\[
 f^m(x)=x+mc.
\]

设 $a$ 是题设给出的不动点。若 $x>a$，由 $f$ 严格增加及 $f(a)=a$，所有正、负整数次迭代 $f^m(x)$ 都应保持大于 $a$。若 $c\ne0$，数列 $x+mc$ 当 $m$ 遍历整数时必有一项小于 $a$，矛盾。因此 $c=0$。对 $x<a$ 同理。于是每个 $x$ 都满足 $T(x)=0$，即
\[
 f(x)\equiv x.
\]
\end{xsolution}

\paragraph{第二组参考题 10}
\begin{xsolution}
只证 $M$ 连续，$m$ 可对 $-f$ 使用同一论证。固定 $x_0\in[a,b]$。由 $f$ 在 $x_0$ 连续，对任意 $\varepsilon>0$ 可取 $\delta>0$，使当 $u,v$ 都在 $x_0$ 的 $\delta$ 邻域内时
\[
 |f(u)-f(v)|<\varepsilon.
\]

先设 $x>x_0$ 且 $x-x_0<\delta$。因为
\[
 M(x)=\max\left\{M(x_0),\max_{x_0\le y\le x}f(y)\right\},
\]
若 $M(x)>M(x_0)$，则最大值由某个 $y\in[x_0,x]$ 取得。此时
\[
 0<M(x)-M(x_0)
 \le f(y)-f(x_0)<\varepsilon,
\]
因为 $M(x_0)\ge f(x_0)$。所以 $|M(x)-M(x_0)|<\varepsilon$。

再设 $x<x_0$ 且 $x_0-x<\delta$。若 $M(x)=M(x_0)$ 无事可证；否则 $M(x_0)$ 的一个最大值点 $y$ 必在 $[x,x_0]$ 中。又 $M(x)\ge f(x)$，故
\[
 0<M(x_0)-M(x)
 \le f(y)-f(x)<\varepsilon.
\]
所以 $M$ 在 $x_0$ 连续。端点只需作单侧论证。故 $M\in C[a,b]$；同理 $m\in C[a,b]$。
\end{xsolution}

\paragraph{第二组参考题 11}
\begin{xsolution}
令
\[
 S=\{x\in[a,b]:f(x)>x\}.
\]
因为 $f(a)>a$ 且 $f$ 单调增加，所以对充分靠近 $a$ 的 $x>a$，有
\[
 f(x)\ge f(a)>x,
\]
故 $S$ 非空且含有大于 $a$ 的点。又因 $f(b)<b$，对充分靠近 $b$ 的 $x<b$，有
\[
 f(x)\le f(b)<x,
\]
所以 $S$ 在 $b$ 附近为空。于是
\[
 c=\sup S
\]
满足 $a<c<b$。

取 $x_n\in S$ 且 $x_n\uparrow c$。由单调性
\[
 f(c)\ge f(x_n)>x_n,
\]
令 $n\to\infty$ 得 $f(c)\ge c$。

另一方面，对每个 $y>c$ 都有 $y\notin S$，故 $f(y)\le y$。取 $y_n\downarrow c$，由单调性
\[
 f(c)\le f(y_n)\le y_n.
\]
令 $n\to\infty$ 得 $f(c)\le c$。因此
\[
 f(c)=c,
\]
且 $c\in(a,b)$。
\end{xsolution}

\paragraph{第二组参考题 12}
\begin{xsolution}
令
\[
 h=f+g.
\]
题设说 $h$ 在 $[0,1]$ 上单调增加，而 $g$ 连续，所以 $f=h-g$。

反设 $f$ 在 $(0,1)$ 中没有零点。由 $f(0)>0$ 与 $g$ 连续、$h$ 单调可知 $f$ 在 $0$ 的右侧某邻域仍为正；同理由 $f(1)<0$ 知 $f$ 在 $1$ 左侧某邻域为负。令
\[
 c=\inf\{x\in[0,1]:f(x)<0\}.
\]
则 $0<c<1$。存在 $x_n\uparrow c$ 使 $f(x_n)>0$，也存在 $y_n\downarrow c$ 使 $f(y_n)<0$。

单调函数 $h$ 的单侧极限存在，并有
\[
 h(c^-)\le h(c)\le h(c^+).
\]
由于 $g$ 连续，减去 $g(c)$ 得
\[
 f(c^-)\le f(c)\le f(c^+).
\]
而由上述两列点及单侧极限，
\[
 f(c^-)\ge0,
 \qquad f(c^+)\le0.
\]
因此只能有
\[
 f(c^-)=f(c)=f(c^+)=0,
\]
这与假定无零点矛盾。所以 $f$ 在 $(0,1)$ 中必有零点。
\end{xsolution}

\paragraph{第二组参考题 13}
\begin{xsolution}
这里“$T_1,T_2$ 不可公约”即 $T_1/T_2\notin\QQ$。反设
\[
 F=f_1+f_2
\]
有正周期 $T$。定义
\[
 h(x)=f_1(x+T)-f_1(x)
 =-[f_2(x+T)-f_2(x)].
\]
从左式可知 $h$ 以 $T_1$ 为周期，从右式可知 $h$ 以 $T_2$ 为周期。由于 $T_1/T_2$ 为无理数，整数线性组合 $mT_1+nT_2$ 在 $\RR$ 中稠密。连续函数 $h$ 具有这一稠密周期群，只能为常值函数，设 $h\equiv c$。

于是
\[
 f_1(x+nT)=f_1(x)+nc.
\]
连续周期函数 $f_1$ 有界，故只能有 $c=0$。所以 $T$ 也是 $f_1$ 的周期。同理 $T$ 也是 $f_2$ 的周期。

非常值连续周期函数若有两个正周期，其比若为无理数，则同样因为整数线性组合形成稠密周期群而必须成为常值函数。因此
\[
 \frac{T}{T_1}\in\QQ,
 \qquad
 \frac{T}{T_2}\in\QQ.
\]
相除得到 $T_1/T_2\in\QQ$，矛盾。故 $f_1+f_2$ 不是周期函数。
\end{xsolution}

\paragraph{第二组参考题 14}
\begin{xsolution}
设 $T>0$ 是 $f$ 的一个周期，$S>0$ 是 $g$ 的一个周期。先证明 $S$ 也是 $f$ 的周期。固定 $x$，利用 $f$ 的周期性，对任意正整数 $n$，
\[
 \begin{aligned}
 |f(x+S)-f(x)|
 &=|f(x+nT+S)-f(x+nT)|\\
 &\le |f(x+nT+S)-g(x+nT+S)|\\
 &\quad+|g(x+nT+S)-g(x+nT)|\\
 &\quad+|g(x+nT)-f(x+nT)|.
 \end{aligned}
\]
中间一项因 $S$ 是 $g$ 的周期而为 $0$，其余两项当 $n\to\infty$ 时都趋于 $0$。故
\[
 f(x+S)=f(x).
\]
所以 $S$ 是 $f$ 的周期。同理 $T$ 也是 $g$ 的周期。

于是 $h=f-g$ 以 $T$ 为周期。任取 $x$，
\[
 h(x)=h(x+nT)\to0
 \qquad(n\to\infty),
\]
所以 $h(x)=0$。故
\[
 f(x)\equiv g(x).
\]
整个证明没有使用连续性。
\end{xsolution}

\paragraph{第二组参考题 15}
\begin{xsolution}
先设 $f$ 一致连续。给定 $\varepsilon>0$，取 $\delta>0$ 使
\[
 |x-y|<\delta\Longrightarrow |f(x)-f(y)|<\varepsilon.
\]
还可由一致连续性对数值 $1$ 取 $\rho>0$，使
\[
 |u-v|<\rho\Longrightarrow |f(u)-f(v)|<1.
\]
把任意 $[x,y]$ 分成长度小于 $\rho/2$ 的若干小段，可得全局估计
\[
 |f(x)-f(y)|
 \le \frac{2}{\rho}|x-y|+1.
\]
因此当 $|x-y|\ge\delta$ 时
\[
 \left|\frac{f(x)-f(y)}{x-y}\right|
 \le\frac2\rho+\frac1\delta.
\]
取
\[
 N>\frac2\rho+\frac1\delta.
\]
若差商绝对值大于 $N$，便只能有 $|x-y|<\delta$，从而 $|f(x)-f(y)|<\varepsilon$。

反过来，设题述差商条件成立。给定 $\varepsilon>0$，取相应的 $N>0$，令
\[
 \delta=\frac\varepsilon{N+1}.
\]
若 $0<|x-y|<\delta$，分两种情况：若
\[
 \left|\frac{f(x)-f(y)}{x-y}\right|>N,
\]
由题设直接有 $|f(x)-f(y)|<\varepsilon$；若差商不超过 $N$，则
\[
 |f(x)-f(y)|\le N|x-y|<N\delta<\varepsilon.
\]
$x=y$ 时更无需讨论。因此 $f$ 一致连续。
\end{xsolution}

\paragraph{第二组参考题 16}
\begin{xsolution}
只讨论每点都是极大值点；每点都是极小值点时对 $-f$ 使用同一证明。

任取有界闭区间 $[c,d]\subset I$。由最值定理，$f$ 在 $[c,d]$ 上取到最小值 $m$，设 $f(x_0)=m$。因为 $x_0$ 又是 $f$ 的局部极大值点，所以存在 $x_0$ 的一个相对邻域 $U$，使
\[
 f(x)\le f(x_0)=m,
 \qquad x\in U.
\]
但 $m$ 是整个 $[c,d]$ 的最小值，又有 $f(x)\ge m$。因此 $f\equiv m$ 于 $U$。

令
\[
 A=\{x\in[c,d]:f(x)=m\}.
\]
由连续性 $A$ 是闭集；上面的论证对 $A$ 中每一点都成立，故 $A$ 在 $[c,d]$ 中也是开集。$A$ 非空，而 $[c,d]$ 连通，所以 $A=[c,d]$。因此 $f$ 在任意闭子区间上恒定，进而在整个开区间 $I$ 上恒定。
\end{xsolution}

\paragraph{第二组参考题 17}
\begin{xsolution}
反设 $f$ 不是常值函数。可取 $a<b$ 使
\[
 f(a)<f(b).
\]
对每个 $c\in(f(a),f(b))$，由介值定理，水平集
\[
 E_c=\{x\in[a,b]:f(x)=c\}
\]
非空。

我们断言每个 $E_c$ 都含有一个非退化区间。令
\[
 A_c=\{x\in[a,b]:f(x)<c\},
 \qquad x_c=\sup A_c.
\]
因为 $f(a)<c$，$A_c$ 非空；又因 $f(b)>c$ 且 $f$ 在 $b$ 连续，$b$ 的某个左邻域中都有 $f>c$，所以 $x_c<b$。由上确界性质可取 $u_n\in A_c$ 使 $u_n\to x_c$。连续性给出 $f(x_c)\le c$。若 $f(x_c)<c$，则由连续性在 $x_c$ 右侧仍有 $f<c$，这与 $x_c$ 是 $A_c$ 的上界矛盾。因此
\[
 f(x_c)=c.
\]
由于有 $u_n\to x_c$ 且 $f(u_n)<c=f(x_c)$，点 $x_c$ 不可能是局部极小值点。题设说每一点都取到极值，所以 $x_c$ 必是局部极大值点。故存在 $\eta>0$，使
\[
 f(x)\le c,
 \qquad x\in(x_c-\eta,x_c+\eta)\cap I.
\]
另一方面，由 $x_c=\sup A_c$，对每个 $x>x_c$ 都不可能有 $f(x)<c$。于是
\[
 f(x)=c,
 \qquad x_c<x<x_c+\eta
\]
（缩小 $\eta$ 使该区间仍在 $I$ 内）。因此 $E_c$ 确实含有非退化区间。

所以，对每个 $c\in(f(a),f(b))$，可在 $E_c$ 中选一个非退化开区间 $J_c$。不同的 $c$ 所对应的 $J_c$ 两两不交。每个 $J_c$ 含一个有理数，于是可把不可列集 $(f(a),f(b))$ 单射到可列集 $\QQ$，矛盾。故 $f$ 只能是常值函数。
\end{xsolution}

\paragraph{第二组参考题 18}
\begin{xsolution}
先讨论严格极大值点。若 $x$ 是严格极大值点，可取有理数 $p_x< x<q_x$，使
\[
 f(t)<f(x),
 \qquad t\in(p_x,q_x),\ t\ne x.
\]
按某个固定规则从所有这样的有理数对中选一对 $(p_x,q_x)$。

若两个不同的严格极大值点 $x\ne y$ 被选到同一有理数对 $(p,q)$，则 $x,y\in(p,q)$。由 $x$ 的严格极大性有 $f(y)<f(x)$，由 $y$ 的严格极大性又有 $f(x)<f(y)$，矛盾。因此
\[
 x\longmapsto(p_x,q_x)
\]
是从严格极大值点集合到 $\QQ^2$ 的单射。$\QQ^2$ 可列，故严格极大值点至多可列。严格极小值点同理。

例如
\[
 f(x)=\cos2\pi x
\]
的所有整数都是严格极大值点，所有半整数 $n+1/2$ 都是严格极小值点，两类点都为可列无穷集。
\end{xsolution}

\paragraph{第二组参考题 19}
\begin{xsolution}
固定 $\varepsilon>0$。在完备度量空间 $[1,2]$ 上，对每个正整数 $N$ 定义闭集
\[
 E_N=\left\{x\in[1,2]:|f(nx)|\le\varepsilon
 \text{ 对所有整数 }n\ge N\right\}.
\]
每个 $E_N$ 是闭集，因为它是闭集
\[
 \{x:|f(nx)|\le\varepsilon\}
\]
的可列交。由题设，对每个固定 $x\in[1,2]$ 都有 $f(nx)\to0$，所以
\[
 [1,2]=\bigcup_{N=1}^\infty E_N.
\]
下面只用闭区间套定理证明：必有某个 $E_N$ 含非退化开区间。若每个 $E_N$ 都没有内点，因为 $E_N$ 闭，则它的补集在任意非退化闭区间中都含有一个非退化开区间。于是可以归纳构造闭区间
\[
 J_1\supset J_2\supset\cdots,
\]
使
\[
 J_N\cap E_N=\varnothing,
 \qquad |J_N|<\frac1N.
\]
由闭区间套定理存在唯一一点 $x_*\in\bigcap_NJ_N$。但对每个 $N$ 都有 $x_*\notin E_N$，这与 $[1,2]=\bigcup_NE_N$ 矛盾。因此某个 $E_N$ 必有内点，从而可取
\[
 [\alpha,\beta]\subset E_N,
 \qquad 1\le\alpha<\beta\le2.
\]
于是对所有整数 $n\ge N$，
\[
 |f(t)|\le\varepsilon,
 \qquad t\in[n\alpha,n\beta].
\]
当 $n$ 足够大时
\[
 (n+1)\alpha\le n\beta,
\]
所以这些区间从某处起彼此重叠，并覆盖一个尾区间 $[M,+\infty)$。因此
\[
 |f(t)|\le\varepsilon,
 \qquad t\ge M.
\]
由于 $\varepsilon>0$ 任意，得到
\[
 \lim_{x\to+\infty}f(x)=0.
\]
\end{xsolution}

\paragraph{第二组参考题 20}
\begin{xsolution}
先证必要性。若 $x_n\to\xi$，则
\[
 x_{n+1}-x_n\to\xi-\xi=0.
\]

现证充分性。设
\[
 x_{n+1}-x_n\to0.
\]
由于 $x_n\in[a,b]$，数列有界。记
\[
 \alpha=\varliminf_{n\to\infty}x_n,
 \qquad
 \beta=\varlimsup_{n\to\infty}x_n.
\]
若 $\alpha=\beta$，数列即收敛。反设 $\alpha<\beta$。

任取 $c\in(\alpha,\beta)$。由下极限和上极限的定义，在任意充分靠后的项中都能找到低于 $c$ 的项，也能找到高于 $c$ 的项。因此可取无穷多次跨越 $c$ 的过程，并在每次过程中取第一次从 $c$ 的一侧到达另一侧的相邻两项 $x_{n_k},x_{n_k+1}$。于是 $c$ 位于这两项之间，故
\[
 |x_{n_k}-c|\le |x_{n_k+1}-x_{n_k}|,
 \qquad
 |x_{n_k+1}-c|\le |x_{n_k+1}-x_{n_k}|.
\]
由 $x_{n+1}-x_n\to0$ 得
\[
 x_{n_k}\to c,
 \qquad x_{n_k+1}\to c.
\]
又 $x_{n_k+1}=f(x_{n_k})$，所以由 $f$ 的连续性
\[
 f(c)=c.
\]
因此 $(\alpha,\beta)$ 中每一点都是 $f$ 的不动点。

现在取 $c=(\alpha+\beta)/2$。由上面的构造，有子列 $x_{n_k}\to c$，故充分大的某一项 $x_N$ 落在 $(\alpha,\beta)$ 内。于是
\[
 x_{N+1}=f(x_N)=x_N,
\]
从而 $x_n=x_N$ 对所有 $n\ge N$ 成立。这使数列收敛，与 $\alpha<\beta$ 矛盾。因此只能有 $\alpha=\beta$，故 $\{x_n\}$ 收敛。
\end{xsolution}
